
% \chapter{Core IOT Platform and Dashboard Design}
\chapter{Task 6 - Data Publishing to CoreIOT Cloud Server}
This chapter explains how the ESP32 publishes telemetry to the CoreIOT (ThingsBoard) MQTT broker, how alarms are generated through rule chains, and how the dashboard renders card values, time-series panels, and alarms. The tone is concise and reproducible.

\section{Device-to-Cloud Path and Broker Setup}

\subsection{End-to-end signal path}
Sensor tasks update a shared \texttt{SensorData\_t} snapshot. The cloud task (\texttt{CORE\_IOT\_task}) serializes that snapshot into MQTT telemetry and pushes it to \texttt{app.coreiot.io}. The ThingsBoard Arduino SDK maintains the MQTT session: telemetry and attributes flow upstream, while RPC commands flow downstream to actuators.

\begin{figure}[H]
    \centering
    % TODO: Replace with actual screenshot of dashboard link in local UI
    \includegraphics[width=0.9\textwidth]{assets/webserver_coreiot_link.png}
    \caption{Local web UI includes a card linking directly to the CoreIOT dashboard.}
    \label{fig:webserver_coreiot_link}
\end{figure}

\subsection{Device provisioning and credentials}
Each device is registered in CoreIOT to obtain its access token (MQTT username). Firmware keeps these parameters in NVS and exposes them in the local web UI so credentials can be updated without reflashing.

\begin{lstlisting}[language=C++, caption={Global configuration variables for Core IOT connection (global.h)}]
// WiFi & CoreIOT Configuration
extern String WIFI_SSID;
extern String WIFI_PASS;
extern String CORE_IOT_TOKEN;
extern String CORE_IOT_SERVER;
extern int CORE_IOT_PORT;
\end{lstlisting}

The connection to CoreIOT uses the ThingsBoard Arduino SDK:

\begin{lstlisting}[language=C++, caption={MQTT client initialization (task\_core\_iot.cpp)}]
constexpr uint32_t MAX_MESSAGE_SIZE = 1024U;

WiFiClient wifiClient;
Arduino_MQTT_Client mqttClient(wifiClient);
ThingsBoard tb(mqttClient, MAX_MESSAGE_SIZE);
\end{lstlisting}

Reconnection authenticates, subscribes RPC/shared attributes, and resends basic attributes:

\begin{lstlisting}[language=C++, caption={Core IOT reconnection and subscription logic}]
void CORE_IOT_reconnect() {
    if (CORE_IOT_TOKEN.isEmpty() || CORE_IOT_SERVER.isEmpty()) return;
    if (!tb.connected()) {
        if (!tb.connect(CORE_IOT_SERVER.c_str(), CORE_IOT_TOKEN.c_str(), CORE_IOT_PORT)) return;

        Serial.println("[CoreIOT] Connected");
        tb.sendAttributeData("macAddress", WiFi.macAddress().c_str());
        tb.RPC_Subscribe(callbacks.cbegin(), callbacks.cend());
        tb.Shared_Attributes_Subscribe(attributes_callback);
        tb.Shared_Attributes_Request(attribute_shared_request_callback);
        tb.sendAttributeData("localIp", WiFi.localIP().toString().c_str());
    } else {
        tb.loop();
    }
}
\end{lstlisting}

To verify the connection and data transmission capabilities, we utilized a Python script to simulate the IoT device behavior before deploying it to the actual hardware. It is critical to note that we specifically selected the \texttt{paho-mqtt} library version 1.6.1 for this implementation, as newer versions (e.g., 2.0) introduce breaking changes in the authentication handling that are incompatible with our current configuration.

\section{Telemetry, RPC, and Reconnect Behavior}

After credentials are set, the cloud task waits for WiFi, connects to MQTT, then cycles telemetry every 10 seconds. RPC keeps cloud switches aligned with actuators, and reconnect logic preserves subscriptions after short outages. Telemetry fields come from a thread-safe snapshot:

\begin{lstlisting}[language=C++, caption={Sensor data structure for telemetry (global.h)}]
typedef struct {
    float temperature;
    float humidity;
    float light_level;
    float moisture_level;
    bool flame_detected;
    bool fan_enabled;
    bool neoled_enabled;
    bool water_pump_enabled;
} SensorData_t;
\end{lstlisting}

The main task runs in a FreeRTOS task, waiting for WiFi connectivity before attempting to connect to Core IOT. It then enters a loop that publishes telemetry data every 10 seconds:

\begin{lstlisting}[language=C++, caption={Telemetry publishing routine in the main task loop}]
void CORE_IOT_task(void *pvParameters) {
    // Wait for WiFi connection
    while (1) {
        if (xSemaphoreTake(xBinarySemaphoreInternet, portMAX_DELAY)) {
            Serial.println("[CoreIOT] WiFi ready");
            break;
        }
        vTaskDelay(500 / portTICK_PERIOD_MS);
    }

    unsigned long lastPublish = 0;
    const unsigned long PUBLISH_INTERVAL = 10000;
    SensorData_t sensorData = {0};

    while (1) {
        CORE_IOT_reconnect();          // MQTT reconnect + tb.loop()
        
        if (millis() - lastPublish >= PUBLISH_INTERVAL) {
            lastPublish = millis();
            getSensorData(&sensorData);  // Thread-safe read
            
            if (tb.connected()) {
                tb.sendTelemetryData("temperature",    sensorData.temperature);
                tb.sendTelemetryData("humidity",       sensorData.humidity);
                tb.sendTelemetryData("light",          sensorData.light_level);
                tb.sendTelemetryData("moisture",       sensorData.moisture_level);
                tb.sendTelemetryData("flame",          sensorData.flame_detected ? 1 : 0);
                tb.sendTelemetryData("fanStatus",      sensorData.fan_enabled ? 1 : 0);
                tb.sendTelemetryData("neoLedEnabled",  sensorData.neoled_enabled ? 1 : 0);
                tb.sendTelemetryData("waterPumpStatus",sensorData.water_pump_enabled ? 1 : 0);
            }
        }
        vTaskDelay(500 / portTICK_PERIOD_MS);
    }
}
\end{lstlisting}

We verify ingestion via the “Latest Telemetry” tab, which shows live updates for all keys above.

\begin{figure}[H]
    \centering
    \includegraphics[width=0.9\textwidth]{assets/lastest_telemetry.png}
    \caption{The “Latest Telemetry” tab confirming live MQTT ingestion.}
    \label{fig:latest_telemetry}
\end{figure}

\textbf{RPC path.} Setter RPCs (\texttt{setFanStatus}, \texttt{setLedEnabled}, \texttt{setWaterPumpStatus}) toggle actuators and refresh the shared snapshot; getter RPCs return the latest state so dashboard switches remain accurate after reconnects. The MQTT reconnect routine resubscribes RPC and shared attributes to avoid drift after short outages.

\textbf{Note on toggle latency.} Switching a control from the CoreIOT dashboard traverses MQTT (cloud broker) before reaching the device, so a brief delay is expected compared to local WebSocket control; the device still updates state and echoes back via RPC getter/telemetry to keep the dashboard in sync.

\section{Rule Chain for Alarms}

Telemetry entering the ThingsBoard Rule Engine is filtered against thresholds: temperature $>$ 33°C or $<$ 15°C, humidity outside 40–70\%, and flame detection. Matching events create alarms (\textit{Environment Warning/Critical/Fire}) and push WebSocket notifications to the dashboard; when values normalize, a clear action resolves the alarm.

\begin{figure}[H]
    \centering
    % TODO: Replace with actual rule chain screenshot
    \includegraphics[width=1\textwidth]{assets/rule_chain_placeholder.png}
    \caption{Rule chain: telemetry filters feed alarm creation and dashboard notifications.}
    \label{fig:rule_chain}
\end{figure}

\section{Dashboard Widgets and Configuration}

\subsection{Overview snapshot}
The dashboard aggregates control switches, status cards, time-series panels, and alarms on one view. The screenshot below shows the live layout with the current session set to a 1-minute window and real-time refresh.

\begin{figure}[H]
    \centering
    % TODO: Replace with the saved screenshot from the dashboard
    \includegraphics[width=1\textwidth]{assets/dashboard_full_placeholder.png}
    \caption{Dashboard overview: device/total cards, RPC switches (pump, fan, Neo LED), light value, WiFi badge, time-series panels, and alarm table.}
    \label{fig:dashboard_full}
\end{figure}

\subsection{Card values and status badges}
Cards display temperature, humidity, soil moisture, light, and WiFi quality. Each card binds to one telemetry key; colors reuse the palette (green/normal, yellow/warning, red/critical).

\begin{figure}[H]
    \centering
    % TODO: Replace with actual dashboard card screenshot
    \includegraphics[width=0.9\textwidth]{assets/dashboard_cards_placeholder.png}
    \caption{Dashboard cards and WiFi badge sourced from live telemetry.}
    \label{fig:dashboard_cards}
\end{figure}

\subsection{Time-series panels}
Dual-axis chart for temperature/humidity and a moisture chart (30\%/70\% guide lines) use a rolling 1-hour window and 5 s refresh. Light can be added as a separate panel if needed. In the captured view, the window is set to 1 minute to highlight real-time stability.

\begin{figure}[H]
    \centering
    % TODO: Replace with actual time-series screenshot
    \includegraphics[width=0.9\textwidth]{assets/timeseries_placeholder.png}
    \caption{Time-series panels for temperature/humidity (dual Y) and soil moisture.}
    \label{fig:timeseries}
\end{figure}

\subsection{Alarm table}
The alarm table lists created time, originator, type, severity, status, and assignee. Severity colors match the palette; clicking an entry opens details and acknowledgement controls.

\begin{figure}[H]
    \centering
    % TODO: Replace with actual alarm table screenshot
    \includegraphics[width=0.9\textwidth]{assets/alarm_table_placeholder.png}
    \caption{Alarm table fed by rule-chain outputs.}
    \label{fig:alarm_table}
\end{figure}

\subsection{Configuration notes}
\begin{itemize}
    \item \textbf{Data keys:} Widget keys must match firmware telemetry names exactly (case-sensitive).
    \item \textbf{RPC names:} Switch widgets call \texttt{setFanStatus}, \texttt{setLedEnabled}, \texttt{setWaterPumpStatus}; paired getters keep toggle state after reconnects.
    \item \textbf{Refresh cadence:} Charts refresh every 5 s to balance responsiveness and load.
    \item \textbf{Color scheme:} Green (normal), yellow (warning), red (critical/fire) is reused across cards, charts, and alarms for quick scanning.
\end{itemize}

